\begin{table}[htbp]
\centering
\caption{问题4.3 实际推广不确定因素与应对策略}
\begin{tabular}{|c|c|c|c|}
\hline
序号 & 不确定因素 & 影响变量 & 应对策略 \\
\hline
1 & 护理人员供给不足与人工成本非线性上涨 & 日固定管理成本、服务能力、响应满意度 & 建立人工成本压力测试，并预留人员补贴或灵活排班方案。 \\
2 & 老人服务需求存在季节性、节假日和突发健康事件波动 & 日服务需求、容量利用率、未满足需求 & 按月或周滚动更新需求预测，保留短期弹性服务资源。 \\
3 & 实际通行时间与表格距离不完全一致 & 距离满意度、覆盖关系、站点归属 & 用步行时间或路网距离替代表格直线距离，并定期校准覆盖矩阵。 \\
4 & 政府补贴标准或单站补贴上限调整 & 政府补贴、利润率、最优定价 & 把补贴金额和补贴上限作为政策参数，定期重算Q3定价。 \\
5 & 服务质量差异影响满意度和复购行为 & 综合满意度、实际有效服务人次、收入 & 引入服务质量评分和投诉率指标，扩展满意度模型。 \\
\hline
\end{tabular}
\end{table}
